\documentclass[11pt]{article}

\usepackage[a4paper, margin=1in]{geometry}
\usepackage{times}
\usepackage{graphicx}
\usepackage{booktabs}
\usepackage{amsmath}
\usepackage{url}
\usepackage{hyperref}
\usepackage{caption}
\usepackage{subcaption}

\title{Analysis of How Machine Learning Research Has Changed Over Time Using arXiv Abstracts}

\author{
\begin{tabular}{c}
Talha Ahmad Khan \\
Sahil Paresh Bhinde \\
Victor Oluwabiyi \\
Chenyu Yang \\
Tian Luo \\
University of Bristol
\end{tabular}
}

\date{}

\begin{document}

\maketitle

\begin{abstract}
% Abstract: around 0.25 pages
% Briefly summarise:
% 1. The overall aim of the project
% 2. The arXiv abstracts dataset
% 3. The main text analytics / machine learning methods used
% 4. The key findings about how ML research has changed over time
\end{abstract}


\section{Introduction}
% Around 1.5 pages
% Include:
% 1. Background: why analysing ML research trends is useful
% 2. Motivation: rapid growth and development of ML research
% 3. Task description: analyse changes in ML research topics over time using arXiv abstracts
% 4. Dataset overview: arXiv abstracts, metadata such as year, title, abstract, categories
% 5. Requirements for a good analysis:
%    - meaningful topic representation
%    - interpretable trends over time
%    - reliable comparison across years
%    - clear visualisation of changes
% 6. Main challenges:
%    - abstracts are short summaries of research
%    - terminology changes over time
%    - topic overlap between ML subfields
%    - possible imbalance across years or categories
% 7. Brief structure of the report


\section{Methods}
% Around 3.5 pages

\subsection{Data Preprocessing and Dataset Preparation}

The dataset for this project is made up of arXiv abstracts from the
\texttt{cs.LG} category, which covers machine learning research. The
papers run from 1999 to 2021, so we have a long enough window to look
at how the field has changed. The original filtered file had 7{,}015
records and five main columns: paper ID, title, abstract, categories,
and publication year. We started by removing one duplicate abstract,
which left us with 7{,}014 papers. We also checked the key columns for
missing values and did not find any.

Next we cleaned the abstract text. Everything was lowercased, extra
whitespace and punctuation were removed, hyphenated terms were split
apart, and the text was tokenised with NLTK. We also removed English
stopwords, since words like ``the'' or ``of'' do not really tell us
anything about the topic of a paper. The cleaned text was kept in a
new \texttt{cleaned\_abstract} column, and the original abstract was
left in place so we could go back to it later if needed. We also added
two simple features for exploration: word count and character count
per abstract.

Because the number of papers per year is very uneven, we did not work
year by year. Instead we grouped the papers into four periods:
\textit{early} (1999--2009), \textit{mid} (2010--2014), \textit{growth}
(2015--2017), and \textit{recent} (2018--2021). The four periods
contain 150, 724, 1{,}488, and 4{,}652 papers, which already shows how
quickly machine learning has grown on arXiv. Most of the dataset comes
from the last few years. Grouping years into periods made the trends
easier to compare, since each bucket has enough papers to be stable.

For modelling we split the data into 70\% train, 15\% validation, and
15\% test. We used stratified sampling on the period label so that each
split keeps the same balance of periods as the full dataset. This gave
us 4{,}909 training, 1{,}052 validation, and 1{,}053 test examples. We
then saved the cleaned dataset, the three splits, and a separate file
per period so the rest of the pipeline could reuse them.

\begin{figure}[h]
    \centering
    \includegraphics[width=\linewidth]{Results/Overview.png}
    \caption{Overview of the arXiv cs.LG dataset, including yearly paper counts, period distribution, and abstract word count distribution.}
    \label{fig:data_overview}
\end{figure}

\subsection{Text Representation and Baseline System}

After preprocessing the dataset, the abstracts were converted into numerical representations so that they could be used for downstream text analytics and machine learning methods. Since the aim of this project is to analyse how machine learning research has changed over time, we compared three different text representation methods: TF-IDF, Word2Vec, and Sentence-BERT. These methods were selected because they capture different aspects of the abstracts, from surface-level vocabulary changes to deeper semantic relationships.

The first representation was TF-IDF, which was used as the baseline representation. TF-IDF represents each abstract as a sparse vector of term importance scores. A term receives a high score if it appears frequently in a specific abstract but is relatively uncommon across the whole corpus. This makes TF-IDF useful for identifying distinctive vocabulary in different time periods. In this project, both unigrams and bigrams were included using an n-gram range of $(1,2)$. This was important because many machine learning concepts are expressed as multi-word phrases, such as \textit{deep learning}, \textit{support vector}, and \textit{reinforcement learning}. English stopwords were removed, and a token pattern was used to keep alphabetic terms with at least three letters, reducing noise from numbers and very short tokens.

In addition to fitting TF-IDF on the full corpus, separate TF-IDF models were also fitted for each time period: \textit{early}, \textit{mid}, \textit{growth}, and \textit{recent}. This allowed us to identify the most distinctive terms within each period. The aim was not only to represent documents for modelling, but also to make vocabulary shifts over time more interpretable. For example, earlier periods were expected to contain more classical machine learning terms, while later periods were expected to show increased use of terms related to neural networks, deep learning, graph-based methods, and other recent research directions.

A count matrix was also created using the same vocabulary settings as the TF-IDF vectorizer. Unlike TF-IDF, the count matrix stores raw word frequencies. This representation was prepared because topic modelling methods such as LDA require raw word counts rather than weighted TF-IDF values. Keeping the same tokenisation and n-gram settings helped make the count-based and TF-IDF-based analyses more comparable.

The second representation was Word2Vec. Unlike TF-IDF, which treats words mainly as independent frequency-based features, Word2Vec learns dense word embeddings based on the surrounding context of each word. Words that appear in similar contexts are placed closer together in the vector space. This is useful for this project because related machine learning concepts may be expressed using different but semantically related terms. The Word2Vec model was trained directly on the cleaned arXiv abstracts, allowing it to learn domain-specific relationships within the machine learning corpus. Each abstract was then represented by averaging the vectors of all words in the abstract. This produced a 100-dimensional vector for each document.

The Word2Vec representation was also checked qualitatively by inspecting similar words for key machine learning terms. For example, terms related to neural networks, reinforcement learning, and Bayesian methods appeared close to related concepts and abbreviations. This suggests that the model captured meaningful relationships within the corpus. However, because the Word2Vec model was trained only on this dataset, it may be less reliable for rare or highly specialised terms.

The third representation was Sentence-BERT. Sentence-BERT encodes the entire abstract as a single dense vector, rather than building a document vector from individual word vectors. We used the \texttt{all-MiniLM-L6-v2} model, which produces 384-dimensional sentence embeddings. This representation was included because it can capture the overall semantic meaning of an abstract more effectively than purely frequency-based methods. This is especially useful when papers from different periods discuss similar research problems using different terminology.

The three representations therefore serve complementary purposes. TF-IDF is the most interpretable and is useful for identifying clear vocabulary changes across time. Word2Vec captures semantic relationships between machine learning concepts and can show connections between related methods. Sentence-BERT provides the richest document-level semantic representation and is useful for identifying conceptually similar papers even when they come from different time periods. Together, these representations allow the project to analyse both surface-level vocabulary shifts and deeper changes in the semantic structure of machine learning research.

The main outputs from this stage were a TF-IDF matrix, a count matrix, Word2Vec document embeddings, and Sentence-BERT document embeddings. These files were saved for use in later comparison, topic analysis, dimensionality reduction, and evaluation stages.

\subsection{Axes of Experimentation}
% Discuss the main experimental directions:
% 1. Time grouping:
%    - yearly analysis
%    - multi-year periods
% 2. Text representation:
%    - Bag-of-Words
%    - TF-IDF
%    - embeddings
% 3. Topic modelling or clustering method:
%    - LDA
%    - K-means
%    - BERTopic, if used
% 4. Dimensionality reduction / visualisation:
%    - PCA
%    - t-SNE
%    - UMAP
% 5. Evaluation and interpretation strategy:
%    - topic coherence
%    - qualitative inspection
%    - trend consistency over time

\subsection{Proposed Improvements}
% Describe each improvement and its hypothesis:
% 1. Improvement 1:
%    e.g. using n-grams to capture phrases such as "deep learning" or "reinforcement learning"
%
% 2. Improvement 2:
%    e.g. using topic modelling to identify latent research themes
%
% 3. Improvement 3:
%    e.g. using embeddings or BERTopic to capture semantic similarity beyond exact keywords
%
% 4. Improvement 4:
%    e.g. grouping years into periods to reduce noise in yearly trends
\subsection{Dimensionality Reduction and Temporal Drift Analysis}

The three text representations described earlier (TF-IDF, Word2Vec,
and Sentence-BERT) all sit in fairly high-dimensional spaces, with
roughly 50{,}000, 100, and 384 dimensions respectively. To see how
machine learning research has shifted between the four periods, and
to get something we can actually plot, we applied PCA to each
representation. We chose PCA because it is simple, deterministic, and
the components come out ordered by how much variance they explain.
That makes it easier to compare results across time and across
representations than something like t-SNE or UMAP, where the output
can change a lot depending on the random seed and the chosen
parameters.

\paragraph{Preparing each representation.}
PCA is sensitive to feature scale, so each representation needed a
small amount of preparation. For Word2Vec and Sentence-BERT, which
produce dense vectors, we mean-centred each feature and scaled it to
unit variance. The mean and standard deviation were computed on the
training set only and then applied to the validation and test sets,
to avoid leaking information between splits. TF-IDF is sparse and very
high-dimensional, so running standard PCA on it directly was not
practical. We used TruncatedSVD instead, which is the usual stand-in
for PCA on sparse matrices. TruncatedSVD does not centre the data,
but our TF-IDF vectors are already L2-normalised, so the components
are still comparable across documents. For simplicity we just call
this step PCA throughout the report.

\paragraph{How many components.}
We used two different numbers of components depending on what we
wanted to do with them. For visualisation we kept only two
components, so that documents could be plotted directly with their
period as a colour. For the quantitative drift analysis we kept 50
components, which gives more stable centroid estimates because more
of the variance is preserved. We also report the cumulative explained
variance for each representation so it is clear how much of the
signal these projections actually keep.

\paragraph{Tracking drift with period centroids.}
To get a number for how much the corpus moves over time, we took the
mean of all training documents in each period within the reduced
space. This gives one centroid per period. Following these centroids
in order (\textit{early} $\rightarrow$ \textit{mid} $\rightarrow$
\textit{growth} $\rightarrow$ \textit{recent}) traces out a path that
shows the overall direction of change. We measured this in two ways.
The first is the Euclidean distance between consecutive centroids,
which tells us how much the average paper changes from one period to
the next. The second is the distance from the \textit{early} centroid
to each later one, which tells us how far things have moved overall.
Because the periods have very different sizes, raw centroids could be
biased toward the bigger periods, so we also recomputed the distances
on bootstrap subsamples of equal size and reported the resulting
confidence intervals.

\paragraph{Checking that the drift is real.}
One thing we worried about is whether the drift we see is actually
about the field changing, or just an artefact of one particular
representation. To check this, we ran the same PCA-based analysis on
all three representations and compared the results. We looked at two
things. First, the order of the four period centroids along the main
components: if the drift is real, the chronological order should
roughly hold across representations. Second, we computed the Spearman
rank correlation between the per-period drift magnitudes from each
pair of representations. When TF-IDF, Word2Vec, and Sentence-BERT
agree, we are more confident that the change is in the research
itself and not in how we encoded the text.

\paragraph{What this stage produces.}
This part of the pipeline feeds three kinds of plots into the results
section. The first is a 2D PCA scatter plot of the whole corpus,
coloured by period, which gives a quick visual sense of how the
papers move over time. The second is a trajectory plot of the four
period centroids in the same 2D space, with arrows between
consecutive periods so the direction of drift is easy to read. The
third is a bar chart of the distances between consecutive periods,
with bootstrap confidence intervals, which gives a number to go with
the picture.

\section{Experimental Evaluation}
% Around 3.5 pages

\subsection{Evaluation Metrics}
% Explain:
% 1. Metrics used for quantitative evaluation
%    - topic coherence, if topic modelling is used
%    - silhouette score, if clustering is used
%    - frequency or proportion changes over time
% 2. Qualitative evaluation
%    - inspecting top words per topic
%    - checking whether topics are meaningful and interpretable
% 3. Limitations of the metrics
%    - topic coherence does not fully capture human interpretability
%    - clustering metrics may not reflect real research fields
%    - frequency trends may be affected by dataset growth

\subsection{Experimental Setup}
% Explain:
% 1. Dataset split or time-based grouping
%    - Since this is trend analysis, the data may be grouped by year rather than train/test split
% 2. How preprocessing was applied
% 3. How models were fitted
% 4. How hyperparameters were selected
% 5. Important libraries or model sources
%    - scikit-learn
%    - gensim
%    - sentence-transformers
%    - BERTopic, if used

\subsection{Results}
% Present:
% 1. Dataset distribution over time
% 2. Top keywords or topics by year / period
% 3. Topic trend plots
% 4. Clustering or embedding visualisations
% 5. Interpretation of how ML research changed over time

\begin{table}[h]
\centering
\begin{tabular}{llp{7cm}}
\toprule
Period / Year & Main Topics & Interpretation \\
\midrule
Early period & -- & -- \\
Middle period & -- & -- \\
Recent period & -- & -- \\
\bottomrule
\end{tabular}
\caption{Summary of major machine learning research themes across different time periods.}
\label{tab:topic_trends}
\end{table}

\subsection{Error Analysis / Limitations of Analysis}
% Since this is not necessarily a classification task, this section can focus on limitations:
% 1. Ambiguous abstracts
% 2. Papers covering multiple topics
% 3. Topic overlap
% 4. Changes in terminology over time
% 5. Dataset bias, e.g. arXiv does not represent all ML research
% 6. Model limitations, e.g. LDA topics may be hard to interpret

\subsection{Discussion}
% Discuss:
% 1. What changes were observed in ML research over time
% 2. Which methods gave the clearest insights
% 3. Strengths and weaknesses of the methods
% 4. Whether findings match known developments in ML
% 5. Recommendations for future analysis


\section{Conclusion}
% Around 1 page
% Summarise:
% 1. Key findings about how ML research changed over time
% 2. Which text analytics methods worked well
% 3. Which methods were less useful
% 4. Main limitations of the analysis
% 5. Open challenges and possible future work


\bibliographystyle{plain}
\bibliography{references}

\end{document}